\scanpage{16}

\mspage{\handwritten{1}}

\begin{flushright}Les Aumettes 15.4.1984\\\end{flushright}
Dear Lipman Bers,\\

Together with Yves Ladegaillerie (a former student of mine) we are running a microseminar \deleted{\unsure{O|o}}on the Teichmüller spaces and groups, my own motivations coming mainly from algebraic geometry, and Ladegaillerie\unsure{(|'}s from his interest in the topology of surfaces. Lately we have met with a problem which I would like to submit to you, as I understand you are the mai\unsure{n} expert on Thurston's hyperbolic geometry approach to Teichmüller space. Before stating the specific problem on hyperbolic ``pants'' (which things boil down to\handwritten{)}, \deleted{L}\handwritten{l}et me tell you what we are really after.

Assuming given a compact oriented surface with boundary $X_0$ \deleted{\unsure{now}} as a reference-surface for constructing the Teichmüller-type spaces, of genus $g$ and with \quad ``holes'' (satisfying $\deleted{\unsure{2}}2g-2+\handwritten{\nu}>0$), my primary interest is in the more ``algebraic'' version of Teichmüller space, corresponding to the question of classifying algebraic non singular curves over $\underline{\mathrm C}$, of genus $g$, with a system of \quad points (all distinct) given on $X$, together with a ``Teichmüller rigidification'' of $(X,S)$ namely a homotopy equivalence between $X_0$ and $X\backslash S$. I'll denote this space, homeomorphic to $\underline{\mathrm C}^{d}$ (where $d=3g-3+\handwritten{\nu}$), by $\overset{\handwritten{\sim}}{\mathrm M}_{\handwritten{g,\nu}}$ (the tilda suggesting that it is the universal covering of a finer object I am still more interested in, namely the algebraic variety (or rather ``multiplicity'', or ``stack'' in the terminology of Mumford-Deligne) of modul\deleted{\unsure{a}}i for algebraic curves of type $(\handwritten{g,\nu})$. Thurston however considers a different modular space, where algebraic curves with a given system of points are replaced by compact conformal oriented surfaces \uline{with boundary}, giving rise to a modular space $\overset{\handwritten{\sim}}{\mathrm{MB}}_{g,\handwritten{\nu}}$ (where the letter B recalls that we are classifying structures with boundary) homeomorphic to $\underline{\mathrm C}^{d}\mathbin{\mathrm x}(\underline{\mathrm R}^{*+})^{\handwritten{\nu}}$, where the extra factor \deleted{$\unsure{\underline{\mathrm C}^{d}}$} corresponds to the extra parameters introducing through the existence of the boundary, namely the length's of the \deleted{various} components of the boundary with respect to the canonical hyperbolic structure on the given surface. Our interest is in pinpointing the precise relationship between the two modular spaces

The obvious idea here is to consider the case of an algebraic curve with \handwritten{$\nu$} points given as a limit-case of a \inserted{compact} conformal surface with boundary, when all the lengths $\deleted{1}\handwritten{\ell}_{\mathrm i}$ \deleted{\unsure{6}}of the components of the boundary tend to zero. Therefore, it looks suitable to consider both modular spaces above as imbedded in a larger third one, which corresponds to the same modular problem as in Thurston's theory, except that we allow the ``bo\deleted{\unsure{p}}undary'' to have some components reduced to just one point, in the neighbourhood of which X is just a conformal surface without boundary, but with a given point
% archive sheet number: 15
